\documentclass[../../main.tex]{subfiles}
\begin{document}
\section{函数列与函数项级数}\label{section:appC-sequences-and-series}

设 $S\subset\mathbb R^n$，假设给定函数 $f:S\to\mathbb R^m$ 与每个整数 $i\geqslant1$ 时的函数 $f_i:S\to\mathbb R^m$。如果对每个 $a\in S$ 和每个 $\varepsilon>0$，存在整数 $N$，使 $i\geqslant N$ 蕴含 $|f_i(a)-f(a)|<\varepsilon$，就称序列 $(f_i)_{i=1}^\infty$ \emph{逐点收敛}到 $f$。如果 $N$ 的选取可与点 $a$ 无关，则称序列\emph{一致收敛}到 $f$：对每个 $\varepsilon>0$ 存在 $N$，使 $i\geqslant N$ 蕴含对所有 $a\in S$ 有 $|f_i(a)-f(a)|<\varepsilon$。如果对每个 $\varepsilon>0$ 存在 $N$，使 $i,j\geqslant N$ 蕴含对所有 $a\in S$ 有 $|f_i(a)-f_j(a)|<\varepsilon$，就称序列\emph{一致 Cauchy}。

\begin{theorem}[一致收敛的性质]\label{theorem:appC-uniform-convergence}
设 $S\subset\mathbb R^n$，假设每个整数 $i\geqslant1$ 时 $f_i:S\to\mathbb R^m$ 连续。
\begin{enumerate}[(a)]
\item 如果 $f_i\to f$ 一致收敛，则 $f$ 连续。
\item 如果序列 $(f_i)_{i=1}^\infty$ 一致 Cauchy，则它一致收敛到某个连续函数。
\item 如果 $f_i\to f$ 一致收敛且 $S$ 是紧的积分区域，则
\begin{align*}
\lim_{i\to\infty}\int_S f_i\,\mathrm{d}V=\int_S f\,\mathrm{d}V.
\end{align*}
\item 如果 $S$ 开，每个 $f_i$ 是 $C^1$ 类，$f_i\to f$ 逐点收敛，且 $(\partial f_i/\partial x^j)$ 当 $i\to\infty$ 时在 $S$ 上一致收敛，则 $\partial f/\partial x^j$ 在 $S$ 上存在且
\begin{align*}
\frac{\partial f}{\partial x^j}=\lim_{i\to\infty}\frac{\partial f_i}{\partial x^j}.
\end{align*}
\end{enumerate}
\end{theorem}

证明见 [Apo74, Rud76, Str00]。

给定 $S\subset\mathbb R^n$ 上一个（实值或向量值）函数的无穷级数 $\sum_{i=0}^\infty f_i$，如果相应部分和序列逐点收敛到某个函数 $f$：
\begin{align*}
f(x)=\lim_{N\to\infty}\sum_{i=0}^N f_i(x)\qquad\text{对所有 }x\in S,
\end{align*}
就称该级数\emph{逐点收敛}。如果其部分和一致收敛，就称级数\emph{一致收敛}。

\begin{proposition}[Weierstrass M 判别法]\label{proposition:appC-weierstrass-m}
假设 $S\subset\mathbb R^n$，$f_i:S\to\mathbb R^k$ 是函数。如果存在正实数 $M_i$，使 $\sup_S|f_i|\leqslant M_i$ 且 $\sum_i M_i$ 收敛，则 $\sum_i f_i$ 在 $S$ 上一致收敛。
\end{proposition}

\begin{exercise}\label{exercise:appC-C-33}
证明命题 C.32。
\end{exercise}
\end{document}
